\section{Introduction}

A model that achieves 95\% accuracy on a benchmark can fail 40\% of the time when challenged by failure-type-specific diagnostics --- but current evaluation frameworks aggregate scores without explaining why. When LLMs fail TruthfulQA questions~\cite{Lin2021TruthfulQA}, we have no automated way to distinguish entity-substitution errors from reasoning failures, leaving correction strategies blind to root causes. For instance, when a model incorrectly answers ``What is the capital of France?'' with ``Lyon'', existing tools report a failure but cannot diagnose whether the model substituted one French city for another (entity-error) or failed to reason about capital-city relationships (reasoning-error). This diagnostic gap prevents targeted correction.

LLM benchmark evaluation serves a critical role in measuring reliability, yet it stops at accuracy scores. State-of-the-art models achieve 60-70\% on TruthfulQA, but the score alone reveals no systematic failure patterns or correction paths. Every incorrect prediction represents a lost opportunity for improvement --- we know \textit{that} the model failed, but not \textit{why}. Without this diagnostic signal, practitioners apply correction methods uniformly: retrieval-augmented generation (RAG) or chain-of-thought (COT) prompting to all failures, regardless of whether the root cause is factual knowledge retrieval or multi-step reasoning.

The deeper problem lies in the disconnection between three independently-evolved research areas. Benchmark evaluation produces aggregate metrics but no individual failure diagnosis. Interpretability research provides insights into model behavior through attention visualization and feature attribution, but these methods require manual example selection, limiting analysis to 10-20 cases per study. Correction methods like RAG and COT have proven effective in various settings, yet they are deployed without failure-type awareness --- a one-size-fits-all approach that misses opportunities to match interventions to root causes.

This gap matters because it prevents scaling interpretability-driven improvement. When failures are not diagnosed, correction strategies cannot target root causes. Uniform RAG application to entity-substitution errors and reasoning failures alike treats symptoms rather than causes. The result: correction effectiveness plateaus, and systematic failure modes remain unaddressed across thousands of benchmark predictions.

We observe that entity-substitution failures exhibit a distinctive attention pattern: when GPT-2 makes entity-substitution errors on TruthfulQA questions, attention entropy over entity spans drops to near-zero (mean = 0.062), indicating the model deterministically attends elsewhere. Non-entity errors show diffuse attention (mean = 0.300). This entropy difference (p < 0.001, Cohen's d = -1.13) provides an actionable classification feature --- threshold-based routing can distinguish failure types with 86.7\% accuracy, enabling matched correction where entity-error routes to RAG (retrieves correct entity for re-attention) rather than uniform application.

Our key insight is that entity-substitution failures are \textit{precision errors} --- the model looks at the wrong place deterministically (60\% exhibit zero entropy, indicating focused-but-wrong attention) rather than \textit{recall errors} (absence of attention). This distinction matters because RAG correction targets entity-level attention misdirection (retrieve correct entity and inject it into context), while COT addresses reasoning-level failures through step-by-step decomposition. Matching the correction method to the failure type targets the root cause.

Building on this insight, we make the following contributions:

\textbf{Methodological.} We demonstrate attention entropy-based failure classification at scale (86.7\% accuracy on 73 processed samples), exceeding the manual analysis baseline of 10-20 examples in prior interpretability work. Our framework automates the pipeline from benchmark failure to attention pattern extraction to diagnostic classification, enabling systematic failure-mode discovery across benchmark-sized datasets.

\textbf{Empirical.} We validate a robust attention pattern signature distinguishing entity-substitution errors from non-entity errors (p = 7.5e-07, Cohen's d = -1.13). Zero-entropy entity-errors (60\% of cases) reveal that entity-substitution failures are precision errors --- the model deterministically ignores the correct entity rather than distributing attention broadly. This observation grounds our matched correction strategy: RAG retrieves the correct entity, targeting the specific attention misdirection.

\textbf{Framework (MOCK VALIDATION ONLY).} We demonstrate matched routing structure (entity-error $\to$ RAG) with synthetic experiments showing +24 percentage point improvement over mismatched routing (entity-error $\to$ COT) using configurable mock success rates (RAG=55$\pm$5\%, COT=30$\pm$5\%). This ablation validates pipeline structure (dual-model framework, gate checking) but \textbf{NOT correction effectiveness} --- real-world deployment with Wikipedia API, GPT-3.5 generation, and GPT-judge evaluation is future work (FW2, HIGH priority). The contribution here is structural feasibility, not deployed effectiveness.

Our framework transforms benchmark evaluation from aggregate scoring to actionable improvement guidance. By connecting individual failures to interpretability-based diagnosis and failure-type-specific correction, we enable practitioners to route entity-errors to RAG without manual analysis. The approach integrates three independently-evolved research areas --- benchmarks, interpretability, and correction --- into a systematic diagnostic workflow.

We organize the paper as follows: Section~\ref{sec:related} discusses related work in benchmark evaluation, interpretability methods, and correction strategies, highlighting the integration gap we address. Section~\ref{sec:methodology} presents our methodology: attention entropy extraction, NER-based span restriction, and threshold-based classification. Section~\ref{sec:experiments} describes our experimental design testing pattern existence (h-e1) and classification utility (h-m1). Section~\ref{sec:results} presents results validating the attention signature and diagnostic routing mechanism. Section~\ref{sec:discussion} discusses zero-entropy entity-errors as precision errors, model-specific pattern limitations, and pending real-world correction validation. Section~\ref{sec:conclusion} concludes with implications for scaling interpretability to automated failure diagnosis.
